This paper argues that stablecoins fundamentally transform parallel exchange rate markets along two key dimensions. First, they reduce transaction frictions and expand access to foreign currency. Second, and more importantly, they generate a publicly observable, high-frequency price that aggregates dispersed information about exchange rate misalignment. As a result, stablecoins act as a public signal embedded in a tradable asset.

A key feature of the environment is that the precision of this public signal is endogenous. As stablecoin usage increases, the stablecoin price becomes more informative about the underlying state of the economy. This creates a feedback loop between asset demand and information aggregation: higher usage improves signal precision, which in turn affects agents' beliefs and behavior.

This mechanism distinguishes the environment from standard global-games models, where the precision of public information is typically taken as exogenous. Here, financial innovation alters not only transaction technology but also the information structure of the economy.
